% =============================================================================
% Section 3: Mechanics and Proof Sketch
% =============================================================================
\section{Mechanics and Proof Sketch}\label{sec:proof}

This section provides the analytical derivation that underlies the
variance decomposition reported in \cref{sec:core-finding}.  The
result is a direct consequence of the linear aggregation rule and the
properties of population covariance.

\subsection{Setup}

Let $A_{ij}$ denote the normalised axis~$j$ score for country~$i$,
$i = 1, \ldots, N$, $j = 1, \ldots, K$, with $N = 27$ and $K = 6$.
The composite is defined as
\begin{equation}\label{eq:composite-def}
  C_i
  \;=\;
  \frac{1}{K} \sum_{j=1}^{K} A_{ij}.
\end{equation}

\subsection{Variance decomposition}

The population variance of the composite across countries is
\begin{align}
  \operatorname{Var}(C)
  &= \operatorname{Var}\!\Bigl(\frac{1}{K}\sum_{j=1}^{K} A_j\Bigr)
  \notag \\[4pt]
  &= \frac{1}{K^2} \sum_{j=1}^{K} \sum_{k=1}^{K}
     \operatorname{Cov}(A_j, A_k)
  \label{eq:var-decomp}
\end{align}
where the random variable notation is understood as operating over
the population of $N$~countries.

Define the \emph{marginal contribution} of axis~$j$ as
\begin{equation}\label{eq:marginal}
  M_j
  \;=\;
  \frac{1}{K^2} \sum_{k=1}^{K} \operatorname{Cov}(A_j, A_k).
\end{equation}

By construction, $\sum_{j=1}^{K} M_j = \operatorname{Var}(C)$, so
each axis's marginal contribution sums exactly to the composite
variance.  The percentage contribution of axis~$j$ is
\begin{equation}\label{eq:pct-contrib}
  \pi_j
  \;=\;
  \frac{M_j}{\operatorname{Var}(C)} \times 100.
\end{equation}

\subsection{Own-variance and cross-covariance decomposition}

Each marginal contribution $M_j$ can be split into two components:
\begin{equation}\label{eq:own-cross}
  M_j
  \;=\;
  \underbrace{\frac{1}{K^2}\operatorname{Var}(A_j)}_{\text{own-variance}}
  \;+\;
  \underbrace{\frac{1}{K^2}\sum_{k \neq j}
    \operatorname{Cov}(A_j, A_k)}_{\text{cross-covariance}}.
\end{equation}

Summing over all axes:
\begin{equation}\label{eq:own-cross-total}
  \operatorname{Var}(C)
  \;=\;
  \underbrace{\frac{1}{K^2}\sum_{j=1}^{K}
    \operatorname{Var}(A_j)}_{\text{total own-variance}}
  \;+\;
  \underbrace{\frac{1}{K^2}\sum_{j=1}^{K}\sum_{k \neq j}
    \operatorname{Cov}(A_j, A_k)}_{\text{total cross-covariance}}.
\end{equation}

\subsection{Interpretation}

Equation~\eqref{eq:var-decomp} is exact: no approximation, no
model, no distributional assumption.  It holds for any linear
aggregation of any set of random variables.  The only structural
requirement is the equal-weight mean in \cref{eq:composite-def}.

The empirical result reported in \cref{sec:core-finding} is a direct
evaluation of \cref{eq:marginal,eq:pct-contrib} on the
$27 \times 6$ matrix of observed axis scores.  No estimation or
inference is involved; $\pi_j$ is a population quantity computed over
the full population of EU-27 countries.

\subsection{Why two axes can dominate}

An axis contributes disproportionately to composite variance when
either or both of the following conditions hold:

\begin{enumerate}[leftmargin=2em]
  \item \textbf{High own-variance.}  The axis scores are more
    dispersed across countries.  For the ISI, the defence axis has
    a population standard deviation of 0.244 and the logistics axis
    0.205, compared to 0.028 for financial.

  \item \textbf{Positive covariance with other axes.}  If high scores
    on axis~$j$ tend to coincide with high scores on other axes, the
    cross-covariance terms in \cref{eq:own-cross} amplify the
    marginal contribution.
\end{enumerate}

\noindent
In the EU-27 data, defence and logistics satisfy both conditions.
Their own-variances are the two highest, and their cross-covariance
terms are predominantly positive (though modest).  The financial
axis fails on both counts: its own-variance is the lowest of the
six, and its correlations with all other axes are near zero
($|t| < 1.15$; see \cref{sec:robustness}).
